\documentclass[10pt,a4paper]{amsart}
\usepackage[margin=19mm]{geometry}
\usepackage{amsmath,amssymb,mathtools}
\usepackage[scr=rsfs,cal=euler]{mathalfa}
\usepackage{xfrac}
\usepackage[unicode,hidelinks,bookmarks=false]{hyperref}
\newcommand{\ie}{i.e.\ }
\newcommand{\cf}{cf.\ }
\newcommand{\resp}{respectively\ }
\newcommand{\Subsection}{\subsection}
\providecommand{\nobreakdash}{\nobreak\hbox{-}}
\setlength{\emergencystretch}{2em}
\multlinegap0pt
\usepackage[shortlabels]{enumitem}
\usepackage{xcolor}
\newtheorem{theorem}{Theorem}
\newtheorem{lemma}{Lemma}
\newcommand{\rev}[1]{#1}
\title{A free boundary problem for spreading of an invasive species in the territory of a native competitor under shifting climate}
\author{Phuoc Vinh Dinh, Phuong Le, Tien Dung Nguyen}
\date{Source: arXiv:2609.14573v1 (CC BY 4.0)}
\begin{document}
\maketitle

\noindent\emph{Source and licence.} A free boundary problem for spreading of an invasive species in the territory of a native competitor under shifting climate. Version arXiv:2609.14573v1 (CC BY 4.0), distributed by arXiv under the Creative Commons Attribution 4.0 International licence, http://creativecommons.org/licenses/by/4.0/, as recorded in the arXiv licence metadata for this exact version and shown on the abstract page https://arxiv.org/abs/2609.14573v1. Full licence notice: \texttt{licenses/arxiv-2609.14573-CC-BY-4.0.txt}.\par\smallskip

\noindent\textit{Editorial scope.} Counted unit: the article's lemma lem:limsup (source0.tex lines 1215--1219). The article's complete original proof of it is delivered with it (source0.tex lines 1221--1336, one \texttt{proof} environment of 116 lines). No mathematics is written, repaired or re-derived: the statement and the proof are the article's own lines, in the article's own notation, and no retained line is substituted or re-flowed.  Retained uncounted as the source-local context the proof needs: lines 97--97; lines 144--153; lines 167--169; lines 316--340; lines 349--431; lines 756--771; lines 759--766.  Nothing was omitted from any retained span, so the package carries no omission record.  No retained line contains a citation, so the wrapper carries no bibliography and no bibliography entry.  No retained line contains a figure environment, an image-inclusion command or drawing code, so no image is delivered and the package declares no asset dependency.  Editorial formatting only, in addition: an AMS \texttt{amsart} preamble that restates the article's own declarations and macros used by the retained lines, namely the environments \texttt{theorem}, \texttt{lemma} on separate counters, together with the standard packages those lines need.  The retained environments print in this document as Theorem 1, Lemma 1 in order of appearance, whereas the article prints the same items with the numbering of its own full document.  The complete native source of the article as distributed in the arXiv e-print for this exact version is bundled unchanged as \texttt{sources/210334/source0.tex}.
	\section{Introduction and statement of the results}\label{sec_1}

	\begin{equation}\label{main}
		\begin{cases}
			u_t = d_1u_{xx} + (A(x - ct) - b_1u - c_1 v) u, & t > 0,~ 0 < x < h(t), \\
			v_t = d_2v_{xx} + (a_2 - b_2u - c_2v)v,          & t > 0,~ 0 < x < +\infty, \\
			u_x(t, 0) = v_x(t, 0) = 0, ~ u(t, x) = 0,        & t > 0, ~ h(t) \leq x < +\infty, \\
			h'(t) = -\mu (A(h(t) - ct))u_x(t, h(t)),          & t > 0, \\
			h(0)=h_0,~ u(0, x) = u_0(x),                       & 0 \leq x \leq h_0, \\
			v(0, x) = v_0(x),                                  & 0 \le x < +\infty,
		\end{cases}
	\end{equation}

	\begin{equation}\label{superior}
		\frac{a_1}{a_2} > \max\left\{\frac{b_1}{b_2}, \frac{c_1}{c_2}\right\}.
	\end{equation}

	\begin{theorem}\label{th:bounds}
		Problem \eqref{main} admits a unique global and uniformly bounded solution $(u,v,h)$ defined for all $t > 0$. Specifically, there exist positive constants $M_1$ and $M_2$ such that
		\[
		0 < u(t,x) \leq M_1
		\quad\text{for }(t, x) \in (0, +\infty) \times [0, h(t)),
		\]
		\[
		0 < v(t,x) \leq M_2
		\quad\text{for }(t, x) \in (0, +\infty) \times [0, +\infty).
		\]
		Moreover, there exists a constant $M_3$ such that
		\[
		0 < h'(t) \leq M_3 \quad\text{for } t \in (0,+\infty).
		\]
		\rev{One may take $M_1=\max\{\|u_0\|_{L^\infty},\,a_1/b_1\}$ and
		$M_2=\max\{\|v_0\|_{L^\infty},\,a_2/c_2\}$; in particular
		\[
		\limsup_{t\to+\infty}u(t,x)\le \frac{a_1}{b_1},\qquad
		\limsup_{t\to+\infty}v(t,x)\le \frac{a_2}{c_2}
		\qquad\text{uniformly for }x\ge0,
		\]
		since $A(\cdot)\le a_1$ and since $u$ and $v$ are subsolutions of the
		corresponding logistic equations.}
		Furthermore, problem \eqref{main} does not have any unbounded solution.
	\end{theorem}

	\begin{theorem}\label{th:comparison}
		Let $(u,v,h)$ be the unique bounded solution of \eqref{main}. Let $T \in (0,+\infty)$ and
		\begin{align*}
			&\underline{h}, \overline{h} \in C^1([0,T]),\\
			&\underline{u} \in C(\overline{U_T}) \cap C^{1,2}(U_T)
			\text{ with } U_T := \{ (t,x) \in \mathbb{R}^2 : t \in (0, T],\, x \in (0, \underline{h}(t))\},\\
			&\overline{u} \in C(\overline{V_T}) \cap C^{1,2}(V_T)
			\text{ with } V_T := \{ (t,x) \in \mathbb{R}^2 : t \in (0, T],\, x \in (0, \overline{h}(t))\},\\
			&\underline{v}, \overline{v}
			\in (L^{\infty} \cap C)([0,T] \times [0,\infty))
			\cap C^{1,2}((0,T] \times [0,\infty)),\\
			&\underline{v}, \overline{v}, \underline{u}, \overline{u} \ge 0,
			~ \underline{h}(0) > 0.
		\end{align*}
		
		Assume that
		\[
		\begin{cases}
			\overline{u}_t \ge d_1\overline{u}_{xx}
			+ (A(x - ct) - b_1\overline{u} - c_1 \underline{v}) \overline{u},
			& 0 < t \le T,~ 0 < x < \overline{h}(t), \\
			\underline{v}_t \le d_2\underline{v}_{xx}
			+ (a_2 - b_2\overline{u} - c_2\underline{v})\underline{v},
			& 0 < t \le T,~ 0 < x < +\infty, \\
			\overline{u}(t, x) = 0,
			& 0 < t \le T, ~ \overline{h}(t) \leq x < +\infty, \\
			\overline{h}'(t) \ge -\mu (A(\overline{h}(t) - ct))\overline{u}_x(t, \overline{h}(t)),
			& 0 < t \le T, \\
			\overline{h}(0) \ge h_0,~ \overline{u}(0, x) \ge u_0(x),
			& 0 \leq x \leq h_0, \\
			\underline{v}(0, x) \le v_0(x),
			& 0 \le x < +\infty,
		\end{cases}
		\]
		and either
		\[
		\overline{u}_x(t, 0) \le 0 \le \underline{v}_x(t, 0) \quad 0 < t \le T
		\]
		or
		\[
		u(t,0) \le \overline{u}(t,0) \text{ and } v(t,0) \ge \underline{v}(t,0)  \quad 0 < t \le T,
		\]
		then
		\[
		h(t) \le \overline{h}(t) \text{ in } (0,T],\quad
		u(t,x) \le \overline{u}(t,x),~ v(t,x) \ge \underline{v}(t,x)
		\text{ for } (t,x) \in (0,T] \times [0,+\infty).
		\]
		
		Similarly, if
		\[
		\begin{cases}
			\underline{u}_t \le d_1\underline{u}_{xx}
			+ (A(x - ct) - b_1\underline{u} - c_1 \overline{v}) \underline{u},
			& 0 < t \le T,~ 0 < x < \underline{h}(t), \\
			\overline{v}_t \ge d_2\overline{v}_{xx}
			+ (a_2 - b_2\underline{u} - c_2\overline{v})\overline{v},
			& 0 < t \le T,~ 0 < x < +\infty, \\
			\underline{u}(t, x) = 0,
			& 0 < t \le T, ~ \underline{h}(t) \leq x < +\infty, \\
			\underline{h}'(t) \le -\mu (A(\underline{h}(t) - ct))\underline{u}_x(t, \underline{h}(t)),
			& 0 < t \le T, \\
			\underline{h}(0) \le h_0,~ \underline{u}(0, x) \le u_0(x),
			& 0 \leq x \leq h_0, \\
			\overline{v}(0, x) \ge v_0(x),
			& 0 \le x < +\infty.
		\end{cases}
		\]		
		and either
		\[
		\underline{u}_x(t, 0) \ge 0 \ge \overline{v}_x(t, 0) \quad 0 < t \le T
		\]
		or
		\[
		u(t,0) \ge \underline{u}(t,0) \text{ and } v(t,0) \le \overline{v}(t,0)  \quad 0 < t \le T,
		\]
		then
		\[
		h(t) \ge \underline{h}(t) \text{ in } (0,T],\quad
		u(t,x) \ge \underline{u}(t,x),~ v(t,x) \le \overline{v}(t,x)
		\text{ for } (t,x) \in (0,T] \times [0,+\infty).
		\]
	\end{theorem}

	\begin{theorem}\label{th:semi-wave} Suppose that $0<c\le c_0$. Then we have the following conclusions:
		\begin{itemize}
			\item[(i)] For any $L\ge0$, the problem
			\begin{equation}\label{semi-wave}
				\begin{cases}
					-c \psi' - d_1 \psi'' = (A(x) - b_1\psi - c_1\varphi)\psi, \quad-\infty < x < L,\\
					-c \varphi' - d_2 \varphi'' = (a_2 - b_2\psi - c_2\varphi)\varphi, \qquad -\infty < x < +\infty,\\
					\psi(-\infty) = \frac{a_1}{b_1}, ~ \psi(x)=0 \text{ for } x\ge L, ~ \psi'(x)<0 \text{ for } x\le L,\\	
					\varphi(-\infty) = 0, ~ \varphi(+\infty)=\frac{a_2}{c_2}, ~ \varphi'(x)>0 \text{ for } x\in\mathbb{R}
				\end{cases}			
			\end{equation}
			has a unique solution $(\psi_L,\varphi_L) \in [C(\mathbb{R})\cap C^2((-\infty, L])] \times C^2(\mathbb{R})$.
			\item[(ii)] The mapping $L \mapsto \psi'_L(L)$ is strictly increasing on $[0, +\infty)$.
			\item[(iii)] There exists a unique $L_0 \geq 0$ such that $-\mu(A(L_0))\psi'_{L_0}(L_0) = c$. Moreover, $L_0 = 0$ if and only if $c = c_0$, and in such a case, we have $(\psi_0,\varphi_0) \equiv (\Psi_{c_0}, \Phi_{c_0})$.
		\end{itemize}
	\end{theorem}

			\begin{equation}
				\begin{cases}
					-c \psi' - d_1 \psi'' = (A(x) - b_1\psi - c_1\varphi)\psi, \quad-\infty < x < L,\\
					-c \varphi' - d_2 \varphi'' = (a_2 - b_2\psi - c_2\varphi)\varphi, \qquad -\infty < x < +\infty,\\
					\psi(-\infty) = \frac{a_1}{b_1}, ~ \psi(x)=0 \text{ for } x\ge L, ~ \psi'(x)<0 \text{ for } x\le L,\\	
					\varphi(-\infty) = 0, ~ \varphi(+\infty)=\frac{a_2}{c_2}, ~ \varphi'(x)>0 \text{ for } x\in\mathbb{R}
				\end{cases}			
			\end{equation}

	\begin{lemma}\label{lem:limsup} Assume $c < c_0$. We have
		\[
		\limsup_{t \to +\infty} \frac{h(t)}{t} \le c.
		\]
	\end{lemma}

	\begin{proof}
		The proof is carried out in two steps.
		
		\textit{Step 1.  A pointwise lower bound for \(v\) at large time.}  
		For any small \(\delta>0\) and any \(T_0>0\), we claim that there exist \(T>T_0\) and \(M>0\) such that
		\[
		v(T,x) \ge \frac{a_2-\delta}{c_2}\quad\text{ for all } x\ge M.
		\]
		
		The argument is standard: we choose $\sigma_0$ such that \(0<\sigma_0<\min\bigl\{\liminf_{x\to\infty}v_0(x),\; a_2/c_2\bigr\}\) and consider the auxiliary problem
		\[
		\begin{cases}
			w_t-d_2w_{xx}=(a_2-c_2w)w, & t>0,\;x>0,\\
			w(t,0)=0,\; w(0,x)=\sigma_0.
		\end{cases}
		\]
		Its solution $w$ converges locally uniformly to the positive stationary solution \(w_*(x)\) of 
		\[
		\begin{cases}
			-d_2w_*''=(a_2-c_2w_*)w_*, & x>0,\\
			w_*(0)=0.
		\end{cases}
		\]
		Moreover, $w'_*(x) > 0$ and $w_*(+\infty) = \frac{a_2}{c_2}$. Hence we can find \(M_1>0\) and \(T>T_0\) such that \(w(t,M_1)\ge (a_2-\delta)/c_2\) for all $t \geq T$.
		Applying the maximum principle to the equation satisfied by $w_x$, we deduce $w_x \ge 0$ for $t >0$ and $x>0$. It follows that
		\begin{equation}\label{limsup1}
			w(t,x) \geq w(t,M_1)\ge \frac{a_2 - \delta}{c_2} \text{ for all } t  \geq T \text{ and } x \geq M_1.
		\end{equation}
		Choose \(M_2\) with \(v_0(x)>\sigma_0\) for \(x\ge M_2\) and set \(M_3=\max\{M_2,h(T)\}\).  Then \(h(t)\le M_3\) for \(t\in[0,T]\) and hence $v$ satisfies
		\[
		v_t - d_2v_{xx} = (a_2 - c_2v)v \text{ for } 0 < t \leq T, ~ x > M_3, \quad v(0,x) > \sigma_0 \text{ for } x > M_3.	
		\]
		
		The function \(\tilde w(t,x)=w(t,x-M_3)\) satisfies \(\tilde w_t-d_2\tilde w_{xx}=(a_2-c_2\tilde w)\tilde w\) for \(0<t\le T\), \(x>M_3\), with \(\tilde w(t,M_3)=0\le v(t,M_3)\) and \(\tilde w(0,x)=\sigma_0<v(0,x)\).  The comparison principle gives
		\begin{equation}\label{limsup2}
			v(t, x) \geq \tilde{w}(t, x) = w(t, x - M_3) \quad \text{ for } 0 < t \leq T, ~ x > M_3.
		\end{equation}
		Taking \(t=T\) and \(M:=M_1+M_3\) and using \eqref{limsup1}, \eqref{limsup2} yield the desired bound.
		
		\textit{Step 2. Construction of an upper solution and conclusion.}
		For each $\delta>0$, we consider the problem
		\begin{equation}\label{limsup3}
			\begin{cases}
				-c \psi' - d_1 \psi'' = (A(x)+2\delta - b_1\psi - c_1\varphi)\psi, ~\quad-\infty < x < L_\delta,\\
				-c \varphi' - d_2 \varphi'' = (a_2 - \delta - b_2\psi - c_2\varphi)\varphi, \quad\qquad -\infty < x < +\infty,\\
				\psi(-\infty) = \frac{a_1 + 2\delta}{b_1}, ~ \psi(x)=0 \text{ for } x\ge L_\delta, ~ \psi'(x)<0 \text{ for } x\le L_\delta,\\	
				\varphi(-\infty) = 0, ~ \varphi(+\infty)=\frac{a_2 - \delta}{c_2}, ~ \varphi'(x)>0 \text{ for } x\in \mathbb{R}.
			\end{cases}
		\end{equation}
		This can be interpreted as a perturbed problem of \eqref{semi-wave}\rev{,
		obtained by replacing $A$ by $A+2\delta$ and $a_2$ by $a_2-\delta$. For
		$\delta>0$ small all the structural assumptions are preserved: $A+2\delta$
		still satisfies the hypotheses on $A$ with $a_0+2\delta<0<a_1+2\delta$,
		condition \eqref{superior} continues to hold by continuity, and $\mu$ is
		defined at $A(L_\delta)+2\delta>a_1$ thanks to the extension of $\mu$
		introduced in Section~\ref{sec_1}}. By Theorem \ref{th:semi-wave}, there exists $c_{0,\delta}>0$ such that if $c \le c_{0,\delta}$ we can find $L_\delta \ge 0$ so that \eqref{limsup3} has a solution \((\psi_\delta,\varphi_\delta)\) with
		\begin{equation}\label{limsup4}
			-\mu\bigl(A(L_\delta)+2\delta\bigr)\,\psi_\delta'(L_\delta)=c.
		\end{equation}
		Moreover, $L_\delta=0$ if and only if $c=c_{0,\delta}$.
		
		Due to the assumption $c < c_0$ and $c_{0,\delta} \to c_0$ as $\delta \to 0$ by continuity, we can fix small $\delta$ such that $c < c_{0,\delta}$. Then we choose $L_\delta \ge 0$ and obtain the solution \((\psi_\delta,\varphi_\delta)\) to problem \eqref{limsup3} with property \eqref{limsup4} as mentioned above.
		
		On the other hand, from the bounds for the original solution \((u,v,h)\) (Theorem \ref{th:bounds}) we can choose \(T_0>0\) so large that
		\[
		u(t,x)\le\frac{a_1+\delta}{b_1}\qquad\text{ for all } t\ge T_0,\;0\le x\le h(t).
		\]
		Applying Step 1 with this $T_0$ and the same $\delta$, we obtain \(T>T_0\) and \(M>0\) with \(v(T,x)\ge (a_2-\delta)/c_2\) for all \(x\ge M\).
		Because \(\psi_\delta(L_\delta)=0\) and \(\psi_\delta\) is decreasing, \rev{and
		since $\psi_\delta(-\infty)=\frac{a_1+2\delta}{b_1}>\frac{a_1+\delta}{b_1}$
		and $\varphi_\delta(-\infty)=0<\inf_{[0,M]}v(T,\cdot)$,} we can pick \(R>0\) so large that
		\[
		L_\delta+R>\max\{h(T),l_0\},\qquad
		\psi_\delta(h(T)\rev{{}-cT}-R)>\frac{a_1+\delta}{b_1},\qquad
		\varphi_\delta(M\rev{{}-cT}-R)<\inf_{0\le x\le M}v(T,x).
		\]
		
		\rev{Now we compare on the time interval $[T,+\infty)$, \emph{keeping the
		original time variable}, which avoids any relabelling of the shifting
		environment. We define, for $t\ge T$,}
		\begin{align*}
			\overline{h}(t)&=c t+L_\delta+R ,\\
			\overline{u}(t,x)&=\psi_\delta\bigl(x-c t-R\bigr) \quad\text{ for } 0\le x \le \overline{h}(t),\\
			\underline{v}(t,x)&=\varphi_\delta\bigl(x-c t-R\bigr)\quad\text{ for } x \ge 0.
		\end{align*}
		We claim that \((\overline{u},\underline{v},\overline{h})\) is an upper solution of the original system \eqref{main} \rev{on $[T,+\infty)$}. Indeed,
		\begin{align*}
			\overline{u}_t-d_1\overline{u}_{xx}
			&=-c\psi_\delta'-d_1\psi_\delta''
			=\bigl(A(x-c t-R)+2\delta-b_1\overline{u}-c_1\underline{v}\bigr)\overline{u} \ge\bigl(A(x-ct)-b_1\overline{u}-c_1\underline{v}\bigr)\overline{u},\\
			\underline{v}_t-d_2\underline{v}_{xx}
			&=-c\varphi_\delta'-d_2\varphi_\delta''
			=\bigl(a_2-\delta-b_2\overline{u}-c_2\underline{v}\bigr)\underline{v} \le\bigl(a_2-b_2\overline{u}-c_2\underline{v}\bigr)\underline{v},\\
			\overline{u}_x(t,0)&=\psi_\delta'( -c t-R)\le0,\\
			\underline{v}_x(t,0)&=\varphi_\delta'( -c t-R)\ge0,\\
			\overline{u}(t,x)&=0 \text{ for } x \ge \overline{h}(t),\\
			\overline{h}'(t)&=c=-\mu\bigl(A(L_\delta)+2\delta\bigr)\,\psi_\delta'(L_\delta)\ge -\mu(A(L_\delta+R))\,\psi_\delta'(L_\delta) = -\mu\bigl(A(\overline{h}(t)-ct)\bigr)\,\overline{u}_x(t,\overline{h}(t)),\\
			\rev{\overline{h}(T)}&\rev{{}=cT+L_\delta+R>h(T)},\\
			\rev{\overline{u}(T,x)}&\rev{{}=\psi_\delta(x-cT-R)\ge\psi_\delta(h(T)-cT-R)>\frac{a_1+\delta}{b_1}\ge u(T,x) \text{ for } 0\le x\le h(T)},\\
			\rev{\underline{v}(T,x)}&\rev{{}=\varphi_\delta(x-cT-R) \le\begin{cases}
				\varphi_\delta(M-cT-R)<\inf_{0\le y\le M}v(T,y)\le v(T,x) & \text{ if } 0\le x\le M,\\
				\varphi_\delta(+\infty)=\frac{a_2-\delta}{c_2}\le v(T,x) & \text{ if } x\ge M.
			\end{cases}}
		\end{align*}
		
		Therefore, $(\overline{u},\underline{v},\overline{h})$ is indeed an upper solution \rev{on $[T,+\infty)$. Note that the proof of Theorem \ref{th:comparison} is insensitive to the choice of the initial time.}  By the comparison principle (Theorem \ref{th:comparison}) we obtain
		\[
		\rev{h(t)\le\overline{h}(t)=ct+L_\delta+R\qquad\forall t\ge T.}
		\]
		Consequently,
		\[
		\limsup_{t\to\infty}\frac{h(t)}{t}\le c.
		\]
		
		This completes the proof of Lemma \ref{lem:limsup}.
	\end{proof}

\end{document}
